\begin{frame}{Példa lefutás - 1. él}
\centering
\begin{tikzpicture}[
  scale=1,
  vtx/.style={circle, draw=black, fill=white, inner sep=1pt, minimum size=16pt},
  mstedge/.style={line width=2.5pt, draw=black},
  otheredge/.style={line width=0.8pt, draw=black},
  font=\small
]

% --- Node coordinates ---
\node[vtx] (a) at (-4.2, 1) {a};
\node[vtx] (b) at (-2.2, 2) {b};
\node[vtx] (d) at (-3, -1.2) {d};
\node[vtx] (e) at (0, 1) {e};
\node[vtx] (f) at (-0.5, -1.2) {f};
\node[vtx] (g) at (2, 2) {g};
\node[vtx] (h) at (2.5, 0) {h};

% -- root node
\node[rootvertex] (c) at (-2.0, 0.2) {c};

% --- MST edges (thick) ---
\draw[mstedge] (d) -- (f)
  node[midway, below=2pt, fill=white, inner sep=1pt, font=\scriptsize]{1};
\draw[mstedge] (g) -- (h)
  node[midway, right=2pt, fill=white, inner sep=1pt, font=\scriptsize]{3};
\draw[mstedge] (e) -- (f)
  node[midway, right=2pt, fill=white, inner sep=1pt, font=\scriptsize]{4};
\draw[mstedge] (b) -- (e)
  node[midway, below=2pt, fill=white, inner sep=1pt, font=\scriptsize]{5};
\draw[mstedge] (c) -- (d)
  node[midway, left=3pt, fill=white, inner sep=1pt, font=\scriptsize]{7};
\draw[mstedge] (a) -- (d)
  node[midway, left=3pt, fill=white, inner sep=1pt, font=\scriptsize]{9};
\draw[mstedge] (e) -- (h)
  node[midway, above=2pt, fill=white, inner sep=1pt, font=\scriptsize]{2};

% --- Non-tree edges (thin) ---
\draw[otheredge] (b) -- (g)
  node[midway, above=2pt, fill=white, inner sep=1pt, font=\scriptsize]{8};
\draw[otheredge] (a) -- (b)
  node[midway, above=3pt, fill=white, inner sep=1pt, font=\scriptsize]{10};
\draw[otheredge] (c) -- (e)
  node[midway, above=2pt, fill=white, inner sep=1pt, font=\scriptsize]{11};
\draw[otheredge] (b) -- (c)
  node[midway, left=2pt, fill=white, inner sep=1pt, font=\scriptsize]{12};
\draw[otheredge] (f) -- (h)
  node[midway, above=2pt, fill=white, inner sep=1pt, font=\scriptsize]{13};
\draw[otheredge, draw=red] (e) -- (g)
  node[midway, above=2pt, fill=white, inner sep=1pt, font=\scriptsize]{6};

% --- DFS (IN, OUT) pairs ---
\node[font=\scriptsize] at ($(a)+(-0.2,0.5)$) {(13, 14)};
\node[font=\scriptsize] at ($(b)+(0,0.6)$) {(9, 10)};
\node[font=\scriptsize] at ($(c)+(-0.6,0.4)$) {(1, 16)};
\node[font=\scriptsize] at ($(d)+(-0.1,-0.6)$) {(2, 15)};
\node[font=\scriptsize] at ($(e)+(0,0.5)$) {(4, 11)};
\node[font=\scriptsize] at ($(f)+(0.0,-0.6)$) {(3, 12)};
\node[font=\scriptsize] at ($(g)+(0,0.6)$) {(6, 7)};
\node[font=\scriptsize] at ($(h)+(0.4,0.5)$) {(5, 8)};

\end{tikzpicture}
\\[0.5em]
\small
\begin{enumerate}
  \item Vertex e is the ancestor of g, then we get the ancestor (ANCESTOR) plan with k1 = IN[e], k2 = IN[g] and thus k1 < k2.
  \item The disjoint sets are linked so g's disjoint set gets h's label. This compresses the subpath (g,h)
\end{enumerate}
\end{frame}

\begin{frame}{Példa lefutás - 1. él (folytatás)}
\centering
\begin{tikzpicture}[
  scale=1,
  vtx/.style={circle, draw=black, fill=white, inner sep=1pt, minimum size=16pt},
  mstedge/.style={line width=2.5pt, draw=black},
  otheredge/.style={line width=0.8pt, draw=black},
  font=\small
]

% --- Node coordinates ---
\node[vtx] (a) at (-4.2, 1) {a};
\node[vtx] (b) at (-2.2, 2) {b};
\node[vtx] (d) at (-3, -1.2) {d};
\node[vtx] (e) at (0, 1) {e};
\node[vtx] (f) at (-0.5, -1.2) {f};
\node[vtx] (g) at (2, 2) {g};
\node[vtx] (h) at (2.5, 0) {h};

% -- root node
\node[rootvertex] (c) at (-2.0, 0.2) {c};

% --- MST edges (thick) ---
\draw[mstedge] (d) -- (f)
  node[midway, below=2pt, fill=white, inner sep=1pt, font=\scriptsize]{1};
\draw[mstedge] (g) -- (h)
  node[midway, right=2pt, fill=white, inner sep=1pt, font=\scriptsize]{3};
\draw[mstedge] (e) -- (f)
  node[midway, right=2pt, fill=white, inner sep=1pt, font=\scriptsize]{4};
\draw[mstedge] (b) -- (e)
  node[midway, below=2pt, fill=white, inner sep=1pt, font=\scriptsize]{5};
\draw[mstedge] (c) -- (d)
  node[midway, left=3pt, fill=white, inner sep=1pt, font=\scriptsize]{7};
\draw[mstedge] (a) -- (d)
  node[midway, left=3pt, fill=white, inner sep=1pt, font=\scriptsize]{9};
\draw[mstedge] (e) -- (h)
  node[midway, above=2pt, fill=white, inner sep=1pt, font=\scriptsize]{2};

% --- Non-tree edges (thin) ---
\draw[otheredge] (b) -- (g)
  node[midway, above=2pt, fill=white, inner sep=1pt, font=\scriptsize]{8};
\draw[otheredge] (a) -- (b)
  node[midway, above=3pt, fill=white, inner sep=1pt, font=\scriptsize]{10};
\draw[otheredge] (c) -- (e)
  node[midway, above=2pt, fill=white, inner sep=1pt, font=\scriptsize]{11};
\draw[otheredge] (b) -- (c)
  node[midway, left=2pt, fill=white, inner sep=1pt, font=\scriptsize]{12};
\draw[otheredge] (f) -- (h)
  node[midway, above=2pt, fill=white, inner sep=1pt, font=\scriptsize]{13};
\draw[otheredge, draw=red] (e) -- (g)
  node[midway, above=2pt, fill=white, inner sep=1pt, font=\scriptsize]{6};

% --- DFS (IN, OUT) pairs ---
\node[font=\scriptsize] at ($(a)+(-0.2,0.5)$) {(13, 14)};
\node[font=\scriptsize] at ($(b)+(0,0.6)$) {(9, 10)};
\node[font=\scriptsize] at ($(c)+(-0.6,0.4)$) {(1, 16)};
\node[font=\scriptsize] at ($(d)+(-0.1,-0.6)$) {(2, 15)};
\node[font=\scriptsize] at ($(e)+(0,0.5)$) {(4, 11)};
\node[font=\scriptsize] at ($(f)+(0.0,-0.6)$) {(3, 12)};
\node[font=\scriptsize] at ($(g)+(0,0.6)$) {(6, 7)};
\node[font=\scriptsize] at ($(h)+(0.4,0.5)$) {(5, 8)};

\end{tikzpicture}
\\[0.5em]
\small
\begin{enumerate}
  \item The next vertex is h since it is the parent of g and then k2 is updated to IN[h].
  \item Continuing the traversal with h (line 12), again line 14 assigns (g,e) to (h,e) where e is the parent of h.
  \item The disjoint sets are linked (line 15) so h's disjoint set gets e’s label. This compresses the subpath (g,h,e).
\end{enumerate}
\end{frame}


\begin{frame}{Példa lefutás - 1. él (folytatás)}
\centering
\begin{tikzpicture}[
  scale=1,
  vtx/.style={circle, draw=black, fill=white, inner sep=1pt, minimum size=16pt},
  mstedge/.style={line width=2.5pt, draw=black},
  otheredge/.style={line width=0.8pt, draw=black},
  font=\small
]

% --- Node coordinates ---
\node[vtx] (a) at (-4.2, 1) {a};
\node[vtx] (b) at (-2.2, 2) {b};
\node[vtx] (d) at (-3, -1.2) {d};
\node[vtx] (e) at (0, 1) {e};
\node[vtx] (f) at (-0.5, -1.2) {f};
\node[vtx] (g) at (2, 2) {g};
\node[vtx] (h) at (2.5, 0) {h};

% -- root node
\node[rootvertex] (c) at (-2.0, 0.2) {c};

% --- MST edges (thick) ---
\draw[mstedge] (d) -- (f)
  node[midway, below=2pt, fill=white, inner sep=1pt, font=\scriptsize]{1};
\draw[mstedge] (g) -- (h)
  node[midway, right=2pt, fill=white, inner sep=1pt, font=\scriptsize]{3};
\draw[mstedge] (e) -- (f)
  node[midway, right=2pt, fill=white, inner sep=1pt, font=\scriptsize]{4};
\draw[mstedge] (b) -- (e)
  node[midway, below=2pt, fill=white, inner sep=1pt, font=\scriptsize]{5};
\draw[mstedge] (c) -- (d)
  node[midway, left=3pt, fill=white, inner sep=1pt, font=\scriptsize]{7};
\draw[mstedge] (a) -- (d)
  node[midway, left=3pt, fill=white, inner sep=1pt, font=\scriptsize]{9};
\draw[mstedge] (e) -- (h)
  node[midway, above=2pt, fill=white, inner sep=1pt, font=\scriptsize]{2};

% --- Non-tree edges (thin) ---
\draw[otheredge] (b) -- (g)
  node[midway, above=2pt, fill=white, inner sep=1pt, font=\scriptsize]{8};
\draw[otheredge] (a) -- (b)
  node[midway, above=3pt, fill=white, inner sep=1pt, font=\scriptsize]{10};
\draw[otheredge] (c) -- (e)
  node[midway, above=2pt, fill=white, inner sep=1pt, font=\scriptsize]{11};
\draw[otheredge] (b) -- (c)
  node[midway, left=2pt, fill=white, inner sep=1pt, font=\scriptsize]{12};
\draw[otheredge] (f) -- (h)
  node[midway, above=2pt, fill=white, inner sep=1pt, font=\scriptsize]{13};
\draw[otheredge, draw=red] (e) -- (g)
  node[midway, above=2pt, fill=white, inner sep=1pt, font=\scriptsize]{6};

% --- DFS (IN, OUT) pairs ---
\node[font=\scriptsize] at ($(a)+(-0.2,0.5)$) {(13, 14)};
\node[font=\scriptsize] at ($(b)+(0,0.6)$) {(9, 10)};
\node[font=\scriptsize] at ($(c)+(-0.6,0.4)$) {(1, 16)};
\node[font=\scriptsize] at ($(d)+(-0.1,-0.6)$) {(2, 15)};
\node[font=\scriptsize] at ($(e)+(0,0.5)$) {(4, 11)};
\node[font=\scriptsize] at ($(f)+(0.0,-0.6)$) {(3, 12)};
\node[font=\scriptsize] at ($(g)+(0,0.6)$) {(6, 7)};
\node[font=\scriptsize] at ($(h)+(0.4,0.5)$) {(5, 8)};

\end{tikzpicture}
\\[0.5em]
\small
\begin{enumerate}
  \item Now the next vertex is e so k2 gets IN[e] making it equal to k1, thus ending the while loop.
  \item The oppositely-oriented edge (e,g) input at line 34 is not processed because eis the ancestor of g and we have already followed the path from descendent to ancestor
\end{enumerate}
\end{frame}

